\documentclass[10pt]{article}
\usepackage[margin=1in]{geometry}
\usepackage[T1]{fontenc}
\usepackage{lmodern,microtype,amsmath,amssymb,amsthm}
\theoremstyle{definition}
\newtheorem{assumption}{Assumption}
\begin{document}
\section*{Stationarity with Gradient-Dependent Hessian Variation}
We consider finding a stationary point of a possibly nonconvex function $f:E\to\mathbb{R}$ on a finite-dimensional real inner product space $E$. The function is three times continuously differentiable, and its Hessian derivative satisfies $\|D^2(\nabla f)(x)\|\leq L_0+L_1\|\nabla f(x)\|$, where $L_0>0$ and $L_1\geq0$. A given point $x^0\in E$ satisfies $f(x^0)-f(x)\leq\Delta$ for every $x\in E$ and $\|\nabla f(x^0)\|\leq G_0$, where $\Delta>0$ and $G_0\geq0$. Separate oracles return the exact gradient and Hessian. For a given $\varepsilon>0$, the output $\bar{x}\in E$ must satisfy $\|\nabla f(\bar{x})\|\leq\varepsilon$ almost surely.

The underlying unconstrained problem is
\[
\inf_{x\in E} f(x),
\]
with success measured by stationarity. The norm on $E$ is induced by its inner product. Write $H(x)=D(\nabla f)(x)$ for the Hessian, viewed as a continuous linear operator from $E$ to $E$; derivatives of operator-valued maps carry the corresponding operator norms. The admissible objectives and oracle access satisfy the following conditions.

\begin{assumption}[Gradient-dependent Hessian variation]
The function $f$ belongs to $C^3(E;\mathbb{R})$, and its Hessian derivative obeys
\[
\|DH(x)\|=\|D^2(\nabla f)(x)\|
\leq L_0+L_1\|\nabla f(x)\|
\qquad \text{for every }x\in E.
\]
\end{assumption}

\begin{assumption}[Initial objective and gradient bounds]
The initial point satisfies
\[
\begin{aligned}
f(x^0)-f(x)&\leq\Delta
&&\text{for every }x\in E,\\
\|\nabla f(x^0)\|&\leq G_0.
\end{aligned}
\]
In particular, $f(x^0)-\Delta$ is a global lower bound on the objective.
\end{assumption}

\begin{assumption}[Exact gradient and Hessian access]
At every call index, every query point $x\in E$, and every realization of oracle randomness, oracle $0$ returns the exact vector $\nabla f(x)$ and oracle $1$ returns the exact operator $H(x)$. Each query costs one call to the selected oracle, with the two call counts recorded separately. Any internal algorithm randomness is independent of the oracle randomness.
\end{assumption}

All parameters $(E,x^0,L_0,L_1,\Delta,G_0,\varepsilon,m)$ are public algorithm inputs, where $m\in\mathbb{N}$ satisfies $m\geq1$. The objective is accessed through the two oracles, and queries may be chosen adaptively anywhere in $E$. Separately specified budgets $K_0,K_1\in\mathbb{N}$ may depend on the full public parameter tuple. The algorithm uses at most $K_0$ gradient calls and $K_1$ Hessian calls, for at most $K_0+K_1$ oracle calls in total.

The algorithm returns a point $\bar{x}\in E$. Writing $\mathbb{P}$ for the joint law of the independent algorithm and oracle randomness, the exact success requirement is
\[
\|\nabla f(\bar{x})\|\leq\varepsilon
\qquad \mathbb{P}\text{-almost surely}.
\]
\end{document}
